% META: {"id":"1SPE-DERLOCAL-CO-021","chapitre":"1SPE-DERIVATION-LOCAL","type_objet":"corrige","exercice_id":"1SPE-DERLOCAL-EX-021","status":"generated"}
% BEGIN-VERIFY
% from sympy import symbols, Eq, solve
% x = symbols('x')
% T = 2*(x - 3) + 2
% assert T.expand() == 2*x - 4
% assert T.subs(x, 3) == 2
% sol = solve(Eq(T, 0), x)
% assert sol == [2]
% END-VERIFY
\begin{corrige}{1SPE-DERLOCAL-EX-021}

\textbf{Question 1.} On calcule :
\[
f(3) = 3^2 - 4 \times 3 + 5 = 9 - 12 + 5 = 2 \qquad \text{et} \qquad f'(3) = 2 \times 3 - 4 = 2.
\]
Donc $\boxed{f(3) = 2}$ et $\boxed{f'(3) = 2}$.

\textbf{Question 2.} L'équation de la tangente au point d'abscisse $a = 3$ est :
\[
y = f'(3)(x - 3) + f(3) = 2(x - 3) + 2 = 2x - 6 + 2 = 2x - 4.
\]
L'équation de la tangente $T$ est $\boxed{y = 2x - 4}$.

\textbf{Question 3.} L'axe des abscisses a pour équation $y = 0$. On résout :
\[
2x - 4 = 0 \iff 2x = 4 \iff x = 2.
\]
La tangente $T$ coupe l'axe des abscisses au point de coordonnées $\boxed{(2\,;\,0)}$.

\textbf{Vérification :} $y = 2 \times 2 - 4 = 0$. \checkmark

\end{corrige}
